% Rozdział aktywny po restrukturyzacji DT.
% Wersja źródłowa: audyt II, jednostka 24

\section{Topologiczne iloczyny tensorowe przestrzeni lokalnie wypukłych}
\label{sec:tensor24}

\providecommand{\atlasref}[2]{#2}


\noindent\textbf{Wymagane wyniki.} Korzystamy z twierdzenia
Hahna--Banacha, definicji nuklearności przez przejścia seminormowe,
rozwinięć Fouriera i Hermite'a. Twierdzenia trwałości nuklearności
dowiedziono w rozdziale~\ref{sec:montel-nuklearne23}; tu korzystamy z nich
w teorii iloczynów tensorowych.
\medskip

W rachunku tensorowym przestrzeni skończenie wymiarowych pojęcie
$E\otimes F$ jest czysto algebraiczne. W analizie funkcjonalnej
potrzebujemy dodatkowo topologii: od niej zależy, które szeregi
tensorów są zbieżne, które formy dwuliniowe są ciągłe i czy
otrzymana przestrzeń ma interpretację jako przestrzeń funkcji.
W tym rozdziale $\K\in\{\R,\C\}$, a przestrzenie lokalnie wypukłe
są Hausdorffa, chyba że zaznaczono inaczej. Zapis $E'$ oznacza
przestrzeń sprzężoną topologicznie z $E$.

\subsection{Produkt tensorowy algebraiczny i własność uniwersalna}

\noindent\emph{Powiązanie tematyczne w Bibliotece matematycznej:} \atlasref{AA-18::aa:tw:tensor-universal}{AA-18: algebraiczna własność uniwersalna iloczynu tensorowego}.


\begin{definicja}[Produkt tensorowy algebraiczny]
\label{def:alg-tensor24}
Niech $E,F$ będą przestrzeniami liniowymi nad $\K$.
Ich \emph{produkt tensorowy algebraiczny} $E\otimes F$ jest
przestrzenią liniową wyposażoną w odwzorowanie dwuliniowe
$\otimes:E\times F\to E\otimes F$, $ (x,y)\mapsto x\otimes y$,
takie, że dla każdej przestrzeni liniowej $G$ i każdego
odwzorowania dwuliniowego $b:E\times F\to G$ istnieje dokładnie
jedno odwzorowanie liniowe $\widetilde b:E\otimes F\to G$
spełniające $\widetilde b(x\otimes y)=b(x,y)$ dla wszystkich
$x\in E$, $y\in F$.
\end{definicja}

\begin{stwierdzenie}[Istnienie i jednoznaczność]
\label{prop:alg-existence24}
Dla dowolnych przestrzeni liniowych $E,F$ produkt tensorowy
algebraiczny istnieje i jest jednoznaczny z dokładnością do
jednoznacznego izomorfizmu zachowującego tensory elementarne.
Każdy jego element jest skończoną sumą $\sum_{j=1}^N x_j\otimes y_j$.
\end{stwierdzenie}
\begin{proof}
Niech $V$ będzie wolną przestrzenią liniową o bazie
$[x,y]$, indeksowanej parami $(x,y)\in E\times F$.
Przez $N\subseteq V$ oznaczmy podprzestrzeń rozpiętą przez
relacje $[x+x',y]-[x,y]-[x',y]$,
$[x,y+y']-[x,y]-[x,y']$, $[ax,y]-a[x,y]$
oraz $[x,ay]-a[x,y]$ dla $a\in\K$.
Przyjmujemy $E\otimes F=V/N$ i $x\otimes y=[x,y]+N$.
Każde $b$ przedłuża się liniowo na $V$ i zeruje się na $N$,
więc przechodzi do jednoznacznego odwzorowania na ilorazie.
To daje własność uniwersalną. W dowolnej realizacji $T$ tensory elementarne rozpinają
całą przestrzeń: odwzorowanie ilorazowe z $T$ do ilorazu
przez ich rozpiętość i odwzorowanie zerowe mają to samo
złożenie z odwzorowaniem dwuliniowym $E\times F\to T$.
Jednoznaczność we własności uniwersalnej oznacza,
że iloraz jest zerowy. Jeśli $T$ i $T'$ są dwiema
realizacjami, własności uniwersalne dają odwzorowania liniowe
w obu kierunkach, których złożenia są identycznościami na
tensorach elementarnych, a więc na całych przestrzeniach.
\end{proof}

\begin{lemat}[Włożenia algebraicznych produktów tensorowych]
\label{lem:alg-inclusion24}
Niech $E,F$ będą przestrzeniami liniowymi nad $\K$,
a $E_0\subseteq E$, $F_0\subseteq F$ ich podprzestrzeniami.
Odwzorowanie kanoniczne $E_0\otimes F_0\to E\otimes F$
jest różnowartościowe. Jeśli $(e_i)_{i=1}^n$ jest bazą
skończenie wymiarowej przestrzeni $E_0$, to każdy element
$E_0\otimes F_0$ ma dokładnie jedno przedstawienie
$\sum_{i=1}^n e_i\otimes y_i$, gdzie $y_i\in F_0$.
\end{lemat}
\begin{proof}
Istnienie przedstawienia wynika z rozpinania przez tensory
elementarne i rozwinięcia pierwszego czynnika w bazie.
Dla funkcjonałów współrzędnych $e_i'$ odwzorowania
$e_i'\otimes I_{F_0}:E_0\otimes F_0\to F_0$,
określone dzięki własności uniwersalnej, odtwarzają
współczynniki $y_i$, co daje jednoznaczność.
W ogólnym przypadku wybieramy algebraiczne dopełnienia
podprzestrzeni $E_0,F_0$ (przez uzupełnienie ich baz).
Odpowiadające im rzuty $P:E\to E_0$ i $Q:F\to F_0$
spełniają $P|_{E_0}=I_{E_0}$, $Q|_{F_0}=I_{F_0}$.
Odwzorowanie $P\otimes Q$ jest lewą odwrotnością
włożenia tensorowego: złożenie jest identycznością na
tensorach elementarnych, więc na całej przestrzeni.
Nie twierdzimy, że rzuty te są ciągłe dla danych topologii.
\end{proof}

\begin{przyklad}[Tensory elementarne a funkcje]
\label{ex:separable-functions24}
Dla zbiorów $X,Y$ oznaczmy przez $\mathcal F(X)$
i $\mathcal F(Y)$ przestrzenie wszystkich funkcji o
wartościach w $\K$. Odwzorowanie
$\mathcal F(X)\otimes\mathcal F(Y)\to\mathcal F(X\times Y)$
określone przez $f\otimes g\mapsto ((x,y)\mapsto f(x)g(y))$
jest liniowe dzięki własności uniwersalnej i różnowartościowe.
Istotnie, zapisujemy tensor w postaci $\sum_{i=1}^n f_i\otimes g_i$
z liniowo niezależnymi $f_i$ (wybierając bazę rozpiętości
pierwszych czynników). Jeśli jego obraz znika, to dla
każdego $y\in Y$ mamy $\sum_i g_i(y)f_i=0$, zatem
wszystkie $g_i(y)=0$. To daje zerowy tensor.
Jeśli $X=Y=\mathbb N$, nie każda funkcja dwóch zmiennych
jest skończoną sumą funkcji rozdzielonych. Na przykład
funkcja $h(i,j)=\delta_{ij}$ nie ma takiego przedstawienia:
dla każdego $N$ jej ograniczenie do $\{1,\ldots,N\}^2$
ma macierz jednostkową rzędu $N$, podczas gdy suma $k$
funkcji rozdzielonych ma macierz rzędu co najwyżej $k$.
W przypadku jednego zbioru skończonego każda funkcja
jest natomiast skończoną sumą takich funkcji. Odpowiednie
topologie i uzupełnienia pozwalają rozważać granice
skończonych sum.
\end{przyklad}

\subsection{Topologia rzutowa}

\begin{definicja}[Seminorma rzutowa]
\label{def:pi-seminorm24}
Niech $E,F$ będą przestrzeniami lokalnie wypukłymi, a $p,q$
ciągłymi seminormami odpowiednio na $E$ i $F$. Dla
$u\in E\otimes F$ definiujemy
\[
 \pi_{p,q}(u)=\inf\Bigl\{\sum_{j=1}^N p(x_j)q(y_j):
 u=\sum_{j=1}^Nx_j\otimes y_j\Bigr\}.
\]
Topologię lokalnie wypukłą generowaną przez wszystkie
$\pi_{p,q}$ nazywamy \emph{topologią rzutową};
odpowiednią przestrzeń oznaczamy $E\otimes_\pi F$,
a jej uzupełnienie $E\widehat\otimes_\pi F$.
\end{definicja}

\begin{lemat}[Własności seminorm rzutowych]
\label{lem:pi-properties24}
Każda funkcja $\pi_{p,q}$ jest seminormą. Dla $x\in E$,
$y\in F$ zachodzi $\pi_{p,q}(x\otimes y)=p(x)q(y)$.
Odwzorowanie kanoniczne $E\times F\to E\otimes_\pi F$
jest ciągłe i dwuliniowe. Jeśli $E,F$ są Hausdorffa,
to również $E\otimes_\pi F$ jest Hausdorffa.
\end{lemat}
\begin{proof}
Jednorodność i nierówność trójkąta wynikają ze skalowania
oraz łączenia przedstawień tensorów. Dla tensora elementarnego
jednoskładnikowe przedstawienie daje nierówność $\le$.
Gdy $p(x)q(y)>0$, twierdzenie Hahna--Banacha daje
$\ell\in E'$ i $m\in F'$ z $|\ell(z)|\le p(z)$,
$|m(w)|\le q(w)$ oraz $\ell(x)=p(x)$, $m(y)=q(y)$
(po odpowiednim wyborze fazy w przypadku zespolonym).
Dla każdego przedstawienia $x\otimes y=\sum_jx_j\otimes y_j$
mamy $p(x)q(y)=|\sum_j\ell(x_j)m(y_j)|
\le\sum_jp(x_j)q(y_j)$. Gdy któryś czynnik zanika,
równość jest natychmiastowa. Ponieważ seminormy
$\pi_{p,q}(x\otimes y)$ są iloczynami $p(x)q(y)$,
ciągłość odwzorowania dwuliniowego w dowolnym punkcie
$(x_0,y_0)$ wynika z oszacowania
\[
\begin{aligned}
\pi_{p,q}(x\otimes y-x_0\otimes y_0)
&\le p(x-x_0)q(y-y_0)\\
&\quad+p(x-x_0)q(y_0)+p(x_0)q(y-y_0).
\end{aligned}
\] Dla $u\ne0$ istnieje para funkcjonałów
liniowych $\ell\in E'$, $m\in F'$ z
$(\ell\otimes m)(u)\ne0$: ograniczamy się do skończenie
wymiarowych przestrzeni rozpiętych przez współczynniki $u$
i przedłużamy funkcjonały współrzędnych twierdzeniem
Hahna--Banacha. Wybieramy $p=|\ell|$, $q=|m|$;
wtedy $\pi_{p,q}(u)\ge|(\ell\otimes m)(u)|>0$.
\end{proof}

\begin{twierdzenie}[Własność uniwersalna topologii rzutowej]
\label{thm:pi-universal24}
Niech $E,F,G$ będą przestrzeniami lokalnie wypukłymi
Hausdorffa. Przyporządkowanie $b\mapsto\widetilde b$ jest bijekcją
między ciągłymi odwzorowaniami dwuliniowymi
$b:E\times F\to G$ i ciągłymi odwzorowaniami liniowymi
$\widetilde b:E\otimes_\pi F\to G$. Jeśli $G$ jest
zupełna, każde takie $\widetilde b$ przedłuża się jednoznacznie
do $E\widehat\otimes_\pi F$.
\end{twierdzenie}
\begin{proof}
Niech $s$ będzie ciągłą seminormą na $G$ i niech $b$
będzie ciągłe. Istnieją zrównoważone otoczenia zera
$U\subset E$, $V\subset F$ takie, że $s(b(x,y))\le1$
na $U\times V$. Po zastąpieniu ich mniejszymi otwartymi
otoczeniami absolutnie wypukłymi i użyciu funkcjonałów
Minkowskiego otrzymujemy ciągłe seminormy $p,q$.
Dla $a>p(x)$ oraz $d>q(y)$, gdzie $a,d>0$,
mamy $x/a\in U$ i $y/d\in V$.
Dwuliniowość daje $s(b(x,y))\le ad$.
Przechodząc do infimum po $a,d$, otrzymujemy
$s(b(x,y))\le p(x)q(y)$, także gdy któraś
z seminorm znika. Stąd dla $u=\sum_jx_j\otimes y_j$ mamy
$s(\widetilde b(u))\le\sum_jp(x_j)q(y_j)$,
a po wzięciu infimum $s(\widetilde b(u))\le\pi_{p,q}(u)$.
To jest ciągłość $\widetilde b$. Odwrotnie, złożenie
ciągłego $\widetilde b$ z ciągłym odwzorowaniem
$(x,y)\mapsto x\otimes y$ jest ciągłe.
Jednoznaczność wynika z rozpinania produktu tensorowego
przez tensory elementarne. Przedłużenie do uzupełnienia
wynika z ciągłości jednostajnej odwzorowania liniowego
oraz zupełności $G$.
\end{proof}

\begin{wniosek}[Forma dwuliniowa i przestrzeń sprzężona]
\label{cor:pi-dual24}
Dla lokalnie wypukłych przestrzeni Hausdorffa $E,F$
przestrzeń liniowa $(E\otimes_\pi F)'$ jest w naturalnej
bijekcji z przestrzenią wszystkich łącznie ciągłych form
dwuliniowych $b:E\times F\to\K$.
\end{wniosek}
\begin{proof}
Stosujemy Twierdzenie~\ref{thm:pi-universal24} dla
$G=\K$. Przyporządkowanie jest liniowe, gdyż
sumowanie i mnożenie przez skalar zachowują
równości na tensorach elementarnych, które
jednoznacznie wyznaczają przedłużenie liniowe.
\end{proof}

\subsection{Topologia injektywna}

\begin{definicja}[Seminorma injektywna]
\label{def:eps-seminorm24}
Niech $E,F$ będą przestrzeniami lokalnie wypukłymi nad $\K$,
a $p,q$ ciągłymi seminormami odpowiednio na $E,F$.
Dla $u\in E\otimes F$ definiujemy
\[
 \varepsilon_{p,q}(u)=
 \sup\Bigl\{\big|\sum_{j=1}^N\ell(x_j)m(y_j)\big|:
 \substack{\ell\in E',\ m\in F',\\
             |\ell(x)|\le p(x)\ (x\in E),\\
             |m(y)|\le q(y)\ (y\in F)}\Bigr\},
 \quad u=\sum_{j=1}^N x_j\otimes y_j.
\]
Topologię wyznaczoną przez te seminormy nazywamy
\emph{topologią injektywną}; piszemy
$E\otimes_\varepsilon F$ i
$E\widehat\otimes_\varepsilon F$ dla jej uzupełnienia.
Występujące w definicji zbiory funkcjonałów są zbiorami
polarnymi domkniętych kul jednostkowych seminorm $p$ i $q$.
Dla przykładu warunek $|\ell(x)|\le1$ przy $p(x)\le1$
pociąga $|\ell(x)|\le p(x)$ przez skalowanie, także gdy
$p(x)=0$ (wtedy wszystkie wielokrotności $x$ leżą w kuli).
Dla ustalonej pary $(\ell,m)$ wartość sumy jest niezależna
od przedstawienia $u$: jest to wartość funkcjonału
$\ell\otimes m$, zadanego przez własność uniwersalną
produktu algebraicznego. Supremum jest skończone, ponieważ
dla dowolnego przedstawienia ogranicza je
$\sum_j p(x_j)q(y_j)$.
\end{definicja}

\begin{twierdzenie}[Porównanie topologii]
\label{thm:eps-pi24}
Dla dowolnych lokalnie wypukłych przestrzeni Hausdorffa
$E,F$ mamy $\varepsilon_{p,q}\le\pi_{p,q}$ dla każdej
pary ciągłych seminorm $p,q$. Obie są seminormami i
$\varepsilon_{p,q}(x\otimes y)=p(x)q(y)$.
Topologia injektywna jest Hausdorffa i nie jest silniejsza
od rzutowej. Kanoniczne odwzorowanie
$E\otimes_\pi F\to E\otimes_\varepsilon F$ jest ciągłe.
\end{twierdzenie}
\begin{proof}
Dla każdego przedstawienia $u=\sum_jx_j\otimes y_j$
i każdej dopuszczalnej pary $(\ell,m)$ zachodzi
$|\sum_j\ell(x_j)m(y_j)|\le\sum_jp(x_j)q(y_j)$.
Bierzemy supremum, a następnie infimum. Własności
seminormy są bezpośrednią konsekwencją definicji,
a równość dla tensora elementarnego wynika z osobnego
zastosowania twierdzenia Hahna--Banacha w obu czynnikach,
tak jak w Lemacie~\ref{lem:pi-properties24}; jeśli
$p(x)q(y)=0$, obie seminormy znikają.
Dla $u\ne0$ wybieramy skończenie wymiarowe podprzestrzenie
$E_0,F_0$ zawierające czynniki pewnego przedstawienia $u$.
Na mocy Lematu~\ref{lem:alg-inclusion24} tensor w
$E_0\otimes F_0$ jest niezerowy. W bazach obu
podprzestrzeni ma więc niezerowy współczynnik,
który wykrywają odpowiednie funkcjonały współrzędnych.
Skończenie wymiarowe podprzestrzenie przestrzeni
lokalnie wypukłych Hausdorffa mają standardową topologię,
więc funkcjonały te są ciągłe i przedłużają się przez
twierdzenie Hahna--Banacha do $\ell\in E'$ i $m\in F'$.
Dla $p=|\ell|$ i $q=|m|$ mamy
$\varepsilon_{p,q}(u)\ge|(\ell\otimes m)(u)|>0$.
To dowodzi oddzielania punktów.
Porównanie seminorm daje ciągłość identyczności z
topologii rzutowej do injektywnej.
\end{proof}

\begin{stwierdzenie}[Funktorialność i symetria iloczynów tensorowych]
\label{prop:dt-tensor-functoriality}
Niech $E,F,E_1,F_1$ będą przestrzeniami lokalnie
wypukłymi Hausdorffa i niech $A:E\to E_1$ oraz
$B:F\to F_1$ będą ciągłymi odwzorowaniami liniowymi.
Odwzorowanie algebraiczne $A\otimes B$, określone przez
$(A\otimes B)(x\otimes y)=Ax\otimes By$, jest ciągłe
zarówno dla topologii rzutowych, jak i injektywnych.
Dla ciągłych seminorm $p_1,q_1$ odpowiednio na $E_1,F_1$
zachodzi, dla $\tau\in\{\pi,\varepsilon\}$,
\[
 \tau_{p_1,q_1}((A\otimes B)u)
 \le\tau_{p_1\circ A,q_1\circ B}(u).
\]
Odwzorowania te przedłużają się jednoznacznie na
odpowiednie uzupełnienia. Ponadto zamiana czynników
$x\otimes y\mapsto y\otimes x$ jest izomorfizmem
liniowo-topologicznym dla obu topologii, również
po uzupełnieniu.
\end{stwierdzenie}
\begin{proof}
Istnienie i jednoznaczność $A\otimes B$ wynikają
z algebraicznej własności uniwersalnej zastosowanej
do odwzorowania dwuliniowego $(x,y)\mapsto Ax\otimes By$.
Oznaczmy $p=p_1\circ A$, $q=q_1\circ B$; są to
ciągłe seminormy. Każde przedstawienie
$u=\sum_j x_j\otimes y_j$ daje przedstawienie obrazu
o sumie iloczynów seminorm równej $\sum_jp(x_j)q(y_j)$.
Wzięcie infimum daje oszacowanie dla $\pi$.
Jeśli $\ell_1,m_1$ są funkcjonałami zdominowanymi
odpowiednio przez $p_1,q_1$, to $\ell_1\circ A$
i $m_1\circ B$ są zdominowane odpowiednio przez $p,q$.
Zatem
\[
 |(\ell_1\otimes m_1)((A\otimes B)u)|
 =|((\ell_1\circ A)\otimes(m_1\circ B))(u)|
 \le\varepsilon_{p,q}(u).
\]
Supremum po dopuszczalnych funkcjonałach daje
oszacowanie dla $\varepsilon$. Oba oszacowania
dowodzą ciągłości. Złożenia oraz identyczności
są zachowywane, ponieważ odpowiadające odwzorowania
zgadzają się na tensorach elementarnych, a te
rozpinają iloczyn algebraiczny.

Ciągłe odwzorowanie liniowe jest jednostajnie ciągłe.
Po złożeniu z kanonicznym włożeniem przestrzeni
docelowej do jej uzupełnienia przedłużamy je
jednoznacznie przez gęstość; uzupełnienie docelowe
jest zupełne i Hausdorffa.
Dla zamiany czynników definicje obu seminorm dają
$\tau_{q,p}(\sum_jy_j\otimes x_j)
=\tau_{p,q}(\sum_jx_j\otimes y_j)$, ponieważ zamiana
jest odwracalna i zamienia odpowiednio przedstawienia
lub pary funkcjonałów. To dowodzi ciągłości w obu
kierunkach. Przedłużenia zamiany i jej odwrotności
są wzajemnie odwrotne przez zgodność na gęstych
iloczynach algebraicznych.
\end{proof}

\begin{twierdzenie}[Dziedziczność normy injektywnej]
\label{thm:eps-injective24}
Niech $X_0\subseteq X$ i $Y_0\subseteq Y$ będą izometrycznie
włożonymi podprzestrzeniami przestrzeni unormowanych.
Dla każdego $u\in X_0\otimes Y_0$ jego norma injektywna
liczona w $X_0\otimes_\varepsilon Y_0$ jest równa normie
liczonej po utożsamieniu z elementem $X\otimes_\varepsilon Y$.
\end{twierdzenie}
\begin{proof}
Utożsamienie tensorów jest uprawnione przez
Lemat~\ref{lem:alg-inclusion24}.
Ograniczenie funkcjonałów z kul jednostkowych $X'$ i $Y'$
daje nierówność $\varepsilon_{X,Y}(u)\le
\varepsilon_{X_0,Y_0}(u)$, gdzie indeksy oznaczają
przestrzenie, w których liczymy normę. Odwrotnie, twierdzenie
Hahna--Banacha przedłuża dowolne $\ell\in X_0'$
i $m\in Y_0'$ z zachowaniem norm. Otrzymujemy więc
identyczny zbiór dopuszczalnych wartości na tensorze $u$.
\end{proof}

\begin{przyklad}[Przypadek skończenie wymiarowy]
\label{ex:finite-dim24}
Dla $E=\K^n$ z dowolną normą oraz przestrzeni Banacha $F$
obydwa uzupełnione produkty tensorowe są liniowo-topologicznie
izomorficzne z $F^n$. Rzeczywiście, przy wyborze bazy
$(e_i)$ i bazy sprzężonej $(e_i')$ każdy tensor ma dokładnie
jedno przedstawienie $u=\sum_{i=1}^ne_i\otimes y_i$.
Z twierdzenia Hahna--Banacha w $F$ mamy
\[
 \|y_i\|=\sup_{\|m\|\le1}|(e_i'\otimes m)(u)|
 \le\|e_i'\|\varepsilon(u).
\]
Z drugiej strony
\[
 \varepsilon(u)\le\pi(u)\le\sum_{i=1}^n\|e_i\|\|y_i\|
 \le\Big(\sum_{i=1}^n\|e_i\|\Big)\max_i\|y_i\|.
\]
Obie normy są zatem równoważne normie maksimum na $F^n$.
Ponieważ $F$ jest zupełna, oba produkty tensorowe są już
zupełne, a ich uzupełnienia nie wprowadzają nowych elementów.
\end{przyklad}

\begin{twierdzenie}[Rozbieżność norm w nieskończonym wymiarze]
\label{thm:hilbert-norms24}
Niech $H$ będzie nieskończenie wymiarową przestrzenią Hilberta
nad $\K$ i $(e_n)$ ortonormalnym ciągiem w $H$. Dla
$u_N=\sum_{n=1}^N e_n\otimes e_n\in H\otimes H$ zachodzi
\[ \varepsilon(u_N)=1,\qquad \pi(u_N)=N. \]
W szczególności na $H\otimes H$ topologie rzutowa i
injektywna są różne.
\end{twierdzenie}
\begin{proof}
Dla $\ell,m$ z kul jednostkowych $H'$ nierówność
Cauchy'ego--Schwarza i Bessela dają
$|\sum_{n=1}^N\ell(e_n)m(e_n)|\le1$. Wybór funkcjonałów
współrzędnych dla $e_1$ pokazuje równość. Z definicji
$\pi(u_N)\le N$. Aby uzyskać nierówność przeciwną,
rozważmy formę dwuliniową
$b(x,y)=\sum_{n=1}^N\langle x,e_n\rangle\langle y,e_n\rangle$.
Przy konwencji liniowości iloczynu skalarnego w pierwszym
argumencie jest ona dwuliniowa. Na mocy Bessela i
nierówności Cauchy'ego--Schwarza ma normę nie większą niż $1$,
a $\widetilde b(u_N)=N$. Dla każdej reprezentacji $u_N$ suma
$\sum_j\|x_j\|\|y_j\|$ jest więc nie mniejsza niż $N$.
Biorąc infimum, dostajemy $\pi(u_N)=N$.
Dla $v_N=u_N/N$ mamy $\varepsilon(v_N)=1/N\to0$,
lecz $\pi(v_N)=1$, co dowodzi różności topologii,
a nie tylko różności wartości norm.
\end{proof}

\subsection{Nuklearność i równość topologii tensorowych}

\begin{lemat}[Seminormy tensorowe i uzupełnienia ilorazów]
\label{lem:tensor-seminorm-completions24}
Niech $E,F$ będą przestrzeniami lokalnie wypukłymi
Hausdorffa, a $p,q$ ich ciągłymi seminormami.
Niech $E_p,F_q$ będą uzupełnieniami unormowanych
ilorazów (Definicja~\ref{def:nuclear-lcs23}).
Dla $u\in E\otimes F$ oznaczmy przez $u_{p,q}$
jego obraz w $E_p\otimes F_q$ pod działaniem
kanonicznych odwzorowań w obu czynnikach. Wówczas
\[
 \pi_{p,q}(u)=\|u_{p,q}\|_\pi,\qquad
 \varepsilon_{p,q}(u)=\|u_{p,q}\|_\varepsilon.
\]
Normy po prawej można liczyć również w odpowiednich
uzupełnionych produktach tensorowych.
\end{lemat}
\begin{proof}
Każde przedstawienie $u=\sum_jx_j\otimes y_j$
daje przedstawienie $u_{p,q}=\sum_j[x_j]_p\otimes[y_j]_q$
o tej samej sumie iloczynów seminorm. Stąd
$\|u_{p,q}\|_\pi\le\pi_{p,q}(u)$.
Jeżeli $\pi_{p,q}(u)=0$, otrzymujemy od razu równość.
W przeciwnym przypadku twierdzenie Hahna--Banacha
dla seminormy $\pi_{p,q}$ daje liniowy funkcjonał
$\ell$ na $E\otimes F$ taki, że
$|\ell(v)|\le\pi_{p,q}(v)$ dla każdego $v$
oraz $\ell(u)=\pi_{p,q}(u)$.
Istotnie, najpierw określamy funkcjonał na
jednowymiarowej podprzestrzeni rozpiętej przez $u$.

Forma $b(x,y)=\ell(x\otimes y)$ spełnia
$|b(x,y)|\le p(x)q(y)$ na mocy
Lematu~\ref{lem:pi-properties24}. Znika, gdy
$x\in\ker p$ lub $y\in\ker q$, więc przechodzi
do formy na ilorazach z normą co najwyżej $1$.
Przedłuża się jednoznacznie do formy dwuliniowej
$\bar b:E_p\times F_q\to\K$ o tym samym oszacowaniu.
Dokładniej, dla ciągów Cauchy'ego $(x_n),(y_n)$
w unormowanych ilorazach różnica wartości formy
jest ograniczona przez
\[
 \|x_n-x_m\|\,\|y_n\|+\|x_m\|\,\|y_n-y_m\|.
\]
Ograniczoność ciągów i to oszacowanie dają istnienie
granicy i jej niezależność od reprezentantów.
Dwuliniowość oraz oszacowanie normowe przechodzą
do granicy. Dla każdego przedstawienia
$u_{p,q}=\sum_j a_j\otimes b_j$ mamy zatem
\[
 \pi_{p,q}(u)=|\widetilde{\bar b}(u_{p,q})|
 \le\sum_j\|a_j\|\,\|b_j\|.
\]
Biorąc infimum, otrzymujemy przeciwną nierówność.
Argument nie zakłada dziedziczności normy rzutowej
dla dowolnych włożeń izometrycznych.

Dla normy injektywnej każdy funkcjonał $\ell$
z $|\ell|\le p$ znika na $\ker p$ i jednoznacznie
przedłuża się do funkcjonału na $E_p$ o normie
co najwyżej $1$. Odwrotnie, każdy taki funkcjonał
na $E_p$ daje przez złożenie z kanonicznym
odwzorowaniem funkcjonał na $E$ zdominowany przez $p$.
To samo dotyczy $q$ i $F_q$. Zbiory dopuszczalnych
wartości w obu definicjach supremum są więc identyczne.
Obie normy tensorowe są normami na iloczynie
przestrzeni unormowanych, a kanoniczne włożenia do
ich uzupełnień są izometriami.
\end{proof}

\begin{twierdzenie}[Nuklearność usuwa różnicę między $\pi$ i $\varepsilon$]
\label{thm:nuclear-pi-eps24}
Niech $E$ będzie nuklearną przestrzenią lokalnie wypukłą
Hausdorffa w sensie Definicji~\ref{def:nuclear-lcs23},
a $F$ dowolną przestrzenią lokalnie wypukłą Hausdorffa.
Wtedy topologie rzutowa i injektywna na
$E\otimes F$ są równe. Ich uzupełnienia są więc
kanonicznie izomorficzne:
$E\widehat\otimes_\pi F\simeq E\widehat\otimes_\varepsilon F$.
\end{twierdzenie}
\begin{proof}
Ustalmy seminormę rzutową $\pi_{p,r}$, gdzie $p$ jest
ciągłą seminormą na $E$, a $r$ na $F$. Z nuklearności
istnieje ciągła seminorma $q\ge p$, taka że kanoniczny
operator $j_{q,p}:E_q\to E_p$ między przestrzeniami
Banacha jest nuklearny. Wybierzmy jego przedstawienie
$j_{q,p}(z)=\sum_{k\ge1}\ell_k(z)a_k$ z
$\sum_k\|\ell_k\|\,\|a_k\|=C<\infty$.
Jeżeli $u=\sum_{j=1}^N x_j\otimes y_j$, to dla każdego
$j$ mamy $[x_j]_p=j_{q,p}([x_j]_q)$.
Wstawiając sumy częściowe przedstawienia nuklearnego
do skończonej sumy tensorów, otrzymujemy w
$E_p\widehat\otimes_\pi F_r$ granicę szeregu
\[
 \sum_{k\ge1}a_k\otimes
 \Bigl(\sum_{j=1}^N\ell_k([x_j]_q)[y_j]_r\Bigr).
\]
Istotnie, dla każdego $j$ sumy częściowe zbiegają
w $E_p$ do $[x_j]_p$, a norma rzutowa tensora
elementarnego jest iloczynem norm.
Z twierdzenia Hahna--Banacha w $F_r$ mamy
\[
 \left\|\sum_j\ell_k([x_j]_q)[y_j]_r\right\|_{F_r}
 \le\|\ell_k\|\,\varepsilon_{q,r}(u),
\]
ponieważ dla $\|m\|_{F_r'}\le1$ funkcjonał
$y\mapsto m([y]_r)$ jest zdominowany przez $r$,
zaś $x\mapsto\ell_k([x]_q)$ jest zdominowany przez
$\|\ell_k\|q$. Dla $\ell_k=0$ nierówność jest oczywista,
a w pozostałych przypadkach dzielimy przez $\|\ell_k\|$.
Z nierówności trójkąta dla normy rzutowej i zbieżności
szeregu otrzymujemy, stosując Lemat~
\ref{lem:tensor-seminorm-completions24},
\[
 \pi_{p,r}(u)=\|u_{p,r}\|_\pi
 \le\sum_k\|a_k\|\,\|\ell_k\|\,\varepsilon_{q,r}(u)
 =C\varepsilon_{q,r}(u).
\]
Otrzymaliśmy, że każda seminorma rzutowa jest
kontrolowana przez seminormę injektywną (z potencjalnie
większą seminormą w pierwszym czynniku). Odwrotna
ciągłość wynika z Twierdzenia~\ref{thm:eps-pi24}.
\end{proof}

\begin{uwaga}[Granica stosowalności]
\label{rem:nuclear-characterization24}
Twierdzenie~\ref{thm:nuclear-pi-eps24} wyjaśnia,
dlaczego dla przestrzeni nuklearnych można niekiedy
pisać $E\widehat\otimes F$ bez wskazywania, czy
użyto topologii $\pi$, czy $\varepsilon$.
Nie oznacza to równości wszystkich możliwych topologii
tensorowych: topologia induktywna $i$, która pojawi się
poniżej, może być silniejsza nawet dla przestrzeni nuklearnych.
Pełne twierdzenie odwrotne Grothendiecka (charakterystyka
nuklearności przez równość $\pi$ i $\varepsilon$ dla wszystkich
czynników) wymaga dodatkowej teorii i nie jest używane w dowodach.
\end{uwaga}

\begin{wniosek}[Przestrzenie Banacha nuklearne tylko w skończonym wymiarze]
\label{cor:banach-finite-nuclear24}
Przestrzeń Banacha $X$ jest nuklearna jako przestrzeń
lokalnie wypukła wtedy i tylko wtedy, gdy ma skończony wymiar.
\end{wniosek}
\begin{proof}
Obie implikacje wykazano w Uwadze~
\ref{rem:nuclear-different23}. Dla pełności przypomnijmy,
że w przypadku nieskończenie wymiarowym nuklearność
wymuszałaby zwartość identyczności na $X$ przez
równoważność normy z dominującą ciągłą seminormą.
W przypadku skończenie wymiarowym dla każdej ciągłej
seminormy $p$ można wybrać $q=p$: identyczność na
skończenie wymiarowym uzupełnieniu $X_p$ jest operatorem
nuklearnym (również gdy $X_p=\{0\}$).
\end{proof}


\begin{stwierdzenie}[Nuklearność produktów]
\label{prop:nuclear-products24}
Niech $(E_\alpha)_{\alpha\in A}$ będzie dowolną rodziną
nuklearnych przestrzeni lokalnie wypukłych Hausdorffa.
Wówczas przestrzeń $E=\prod_{\alpha\in A}E_\alpha$
z topologią produktu jest nuklearna.
\end{stwierdzenie}
\begin{proof}
Niech $p$ będzie ciągłą seminormą na $E$.
Z definicji topologii produktu istnieje skończony zbiór
$J\subseteq A$ taki, że $p(x)=0$, gdy $x_\alpha=0$
dla każdego $\alpha\in J$. Istotnie, w otoczeniu zera
$\{x:p(x)<1\}$ mieści się otoczenie cylindryczne
wyznaczone przez skończenie wiele współrzędnych.
Wektor zerujący się na nich można pomnożyć przez dowolny
skalar bez opuszczania cylindra, więc jego seminorma
musi być zerowa. Jeżeli $J=\varnothing$, to $p=0$ i wystarcza $q=0$.
W przeciwnym przypadku niech $\iota_\alpha:E_\alpha\to E$
będzie włożeniem na współrzędną $\alpha$, z zerami
na pozostałych współrzędnych. Jest ono ciągłe.
Dla $p_\alpha=p\circ\iota_\alpha$ mamy
\[
 p(x)=p\Bigl(\sum_{\alpha\in J}\iota_\alpha(x_\alpha)\Bigr)
 \le\sum_{\alpha\in J}p_\alpha(x_\alpha)=P(x),
\]
ponieważ różnica wektorów należy do $\ker p$.
W szczególności wszystkie $p_\alpha$ są ciągłe.
Dla każdego $\alpha\in J$ wybieramy $q_\alpha\ge p_\alpha$
takie, że $j_{q_\alpha,p_\alpha}$ jest nuklearny,
i przyjmujemy $q(x)=\sum_{\alpha\in J}q_\alpha(x_\alpha)$.
Dla $t=q$ oraz $t=P$ jądro seminormy składa się
z wektorów, których współrzędne z $J$ leżą w jądrach
odpowiednich seminorm składowych. Uzupełnienie $E_t$
identyfikujemy izometrycznie ze skończoną sumą
$\ell^1$ odpowiednich uzupełnień współrzędnych:
odwzorowanie ilorazu na tę sumę zachowuje normę
i ma gęsty obraz przez gęstość każdego składnika.
Przejście $E_q\to E_P$ jest więc sumą
\[
 j_{q,P}=\sum_{\alpha\in J}
 I_\alpha j_{q_\alpha,p_\alpha}R_\alpha,
\]
gdzie $R_\alpha$ jest rzutem współrzędnym, a $I_\alpha$
włożeniem do sumy $\ell^1$; ich normy nie przekraczają $1$.
Zgodność sprawdzamy na gęstym obrazie $E$.
Każdy składnik jest nuklearny na mocy
Stwierdzenia~\ref{prop:nuclear-ideal23}, zatem
skończona suma też jest nuklearna. Przez $p\le P$
operator $j_{P,p}$ jest ograniczony, a złożenie
$j_{q,p}=j_{P,p}j_{q,P}$ jest nuklearne.
Ponadto $q\ge p$, co kończy dowód.
\end{proof}

\begin{uwaga}[Trwałość względem podprzestrzeni i granic induktywnych]
\label{rem:nuclear-permanence24}
Trwałość nuklearności względem podprzestrzeni i ilorazów
została udowodniona w
Twierdzeniu~\ref{thm:nuclear-subquot23}, a względem
przeliczalnych granic induktywnych --- w
Twierdzeniu~\ref{thm:nuclear-indlim23}.
Dlatego nuklearność $\mathcal D_K(\Omega)$ i
$\mathcal D(\Omega)$ jest już wynikiem
Twierdzenia~\ref{thm:nuclear-D23}, a nie osobnym
nieudowodnionym założeniem. Wyniki te nie uprawniają
do utożsamienia topologii $i$ i $\pi$ na
$\mathcal D(\Omega)\otimes\mathcal D(\Lambda)$:
porównują one odmienne własności uniwersalne.
\end{uwaga}

\subsection{Przykłady produktów tensorowych przestrzeni funkcyjnych}

\begin{twierdzenie}[Produkt tensorowy $\ell^1$]
\label{thm:l1-tensor24}
Niech $F$ będzie przestrzenią Banacha. Wówczas istnieje
kanoniczny izometryczny izomorfizm
$\ell^1(\N)\widehat\otimes_\pi F\simeq\ell^1(\N;F)$,
gdzie prawa strona składa się z ciągów $(y_n)$ spełniających
$\sum_n\|y_n\|_F<\infty$, z normą tej sumy.
\end{twierdzenie}
\begin{proof}
Najpierw definiujemy liniowe odwzorowanie
$J:\ell^1\otimes F\to\ell^1(\N;F)$ wzorem
\[
 J\Bigl(\sum_{j=1}^M a^{(j)}\otimes y_j\Bigr)
 =\Bigl(\sum_{j=1}^M a_n^{(j)}y_j\Bigr)_{n\ge1}.
\]
Jest ono dobrze określone dzięki własności uniwersalnej
produktu algebraicznego: odwzorowanie
$(a,y)\mapsto(a_ny)_{n\ge1}$ jest dwuliniowe i jego
wartości należą do $\ell^1(\N;F)$, ponieważ
$\sum_n\|a_n y\|=\|a\|_1\|y\|$.
Przestrzeń $\ell^1(\N;F)$ jest zupełna: dla ciągu
Cauchy'ego granice współrzędnych istnieją w $F$,
a oszacowania norm różnic przechodzą do granicy
najpierw dla skończonych sum, następnie przez
supremum po długościach tych sum. Daje to
sumowalność norm współrzędnych granicy i zbieżność w $\ell^1$.
Dla dowolnego przedstawienia
$u=\sum_j a^{(j)}\otimes y_j$ zachodzi
$\|Ju\|_{\ell^1(F)}\le\sum_j\|a^{(j)}\|_1\|y_j\|$,
a więc $\|Ju\|_{\ell^1(F)}\le\pi(u)$.

Jeżeli $u=\sum_{n=1}^N e_n\otimes y_n$, to nierówność
przeciwna wynika z twierdzenia Hahna--Banacha.
Wybieramy $m_n\in F'$ z $\|m_n\|\le1$ i
$m_n(y_n)=\|y_n\|$ (dla $y_n=0$ bierzemy $m_n=0$).
Forma $b(a,y)=\sum_{n=1}^Na_nm_n(y)$ ma normę
nie większą niż $1$ na $\ell^1\times F$ oraz
$\widetilde b(u)=\sum_n\|y_n\|$. Zatem
$\pi(u)=\sum_n\|y_n\|=\|Ju\|_{\ell^1(F)}$
dla tensorów o skończonym nośniku współrzędnych.

Pozostaje uzasadnić, że wyliczenie normy obejmuje
\emph{dowolny} tensor algebraiczny. Jeśli
$u=\sum_{j=1}^M a^{(j)}\otimes y_j$, niech $P_N$
obcina ciąg po $N$ współrzędnych i połóżmy
$u_N=\sum_j(P_Na^{(j)})\otimes y_j$. Wtedy
\[
 \pi(u-u_N)\le\sum_{j=1}^M
 \|a^{(j)}-P_Na^{(j)}\|_1\,\|y_j\|\longrightarrow0.
\]
Ponadto $\|Ju_N-Ju\|_{\ell^1(F)}\le\pi(u_N-u)\to0$. Równość
norm dla $u_N$, wraz z ciągłością norm obu przestrzeni,
daje $\pi(u)=\|Ju\|_{\ell^1(F)}$.
Odwzorowanie $J$ jest więc izometrią o obrazie
zawierającym wszystkie ciągi o skończonym nośniku.
Przedłuża się do izometrii z zupełnego
$\ell^1\widehat\otimes_\pi F$ na $\ell^1(\N;F)$:
jej obraz jest domknięty i zarazem gęsty.
\end{proof}

\begin{twierdzenie}[Jądra gładkie i szybko malejące]
\label{thm:smooth-schwartz-tensor24}
Dla niepustych zbiorów otwartych $\Omega\subseteq\R^{d_1}$,
$\Lambda\subseteq\R^{d_2}$ oraz $d_1,d_2\ge1$ mamy
naturalne izomorfizmy liniowo-topologiczne
\[
 \mathcal E(\Omega)\widehat\otimes_\pi\mathcal E(\Lambda)
 \simeq \mathcal E(\Omega\times\Lambda),\qquad
 \mathcal S(\R^{d_1})\widehat\otimes_\pi\mathcal S(\R^{d_2})
 \simeq \mathcal S(\R^{d_1+d_2}).
\]
W obu wzorach można zastąpić $\pi$ przez $\varepsilon$.
\end{twierdzenie}
\begin{proof}
W obu przypadkach odwzorowanie kanoniczne na tensorach
elementarnych przypisuje $f\otimes g$ funkcję
$(x,y)\mapsto f(x)g(y)$.

Dla przestrzeni Schwartza stosujemy izomorfizmy z
przestrzeniami szybko malejących ciągów, udowodnione
w Lemacie~\ref{lem:hermite-schwartz23}.
Izomorfizmy w czynnikach indukują izomorfizm
uzupełnionych iloczynów rzutowych na mocy
Stwierdzenia~\ref{prop:dt-tensor-functoriality},
zastosowanego do obu izomorfizmów oraz ich odwrotności.
Z twierdzenia Fubiniego współczynniki Hermite'a
funkcji $(x,y)\mapsto f(x)g(y)$ są iloczynami
współczynników obu funkcji (całki są bezwzględnie
zbieżne dzięki nierówności Cauchy'ego--Schwarza).
Wystarczy więc zbadać odwzorowanie
$J:s(I)\otimes s(J_0)\to s(I\times J_0)$,
gdzie $I=\N_0^{d_1}$, $J_0=\N_0^{d_2}$,
a tensor elementarny $(a_i)\otimes(b_j)$
przechodzi na macierz $(a_i b_j)_{i,j}$.
Oznaczmy $w(i,j)=1+|i|+|j|$ oraz
$P_m(c)=\sup_{i,j}w(i,j)^m|c_{ij}|$.
Ponieważ $w(i,j)\le(1+|i|)(1+|j|)$,
dla każdego przedstawienia tensora $u$ otrzymujemy
$P_m(Ju)\le\sum_k p_m(a^{(k)})p_m(b^{(k)})$.
Biorąc infimum, mamy
$P_m(Ju)\le\pi_{p_m,p_m}(u)$.
Seminormy o równych indeksach są współkońcowe
wśród par seminorm obu przestrzeni, ponieważ
$p_k$ rosną z $k$. Zupełność $s(I\times J_0)$
(Twierdzenie~\ref{thm:nuclear-s23}) daje ciągłe
przedłużenie $\widehat J$ do uzupełnienia.

Dla $c\in s(I\times J_0)$ definiujemy odwzorowanie
odwrotne $R(c)$ jako sumę szeregu
$\sum_{i,j}c_{ij}e_i\otimes e_j$ w uzupełnieniu.
Dla $M>2m+d_1+d_2$ mamy
\[
 \sum_{i,j}|c_{ij}|(1+|i|)^m(1+|j|)^m
 \le P_M(c)\sum_{i,j}w(i,j)^{2m-M}<\infty.
\]
Ostatnia suma jest skończona na mocy zliczania
wielowskaźników w Twierdzeniu~\ref{thm:nuclear-s23}
dla wymiaru $d_1+d_2$. Sumy po skończonych zbiorach
indeksów tworzą sieć Cauchy'ego w każdej seminormie
rzutowej; z zupełności istnieje granica.
Powyższe oszacowanie dowodzi ciągłości $R$.
Współczynniki $\widehat J R(c)$ są równe $c_{ij}$,
ponieważ każda ewaluacja współrzędnej jest ciągła.
Dla tensora elementarnego $a\otimes b$ sumy po
prostokątnych zbiorach $|i|\le N$, $|j|\le N$
są tensorami obcięć $a^{[N]}\otimes b^{[N]}$.
Obcięcia zbiegają do $a,b$ we wszystkich normach
$p_m$, gdyż ogony są ograniczone przez
$p_{m+1}(a)/(1+N)$ i $p_{m+1}(b)/(1+N)$.
Z ciągłości kanonicznego odwzorowania dwuliniowego
wynika $R J(a\otimes b)=a\otimes b$.
Liniowość daje równość na iloczynie algebraicznym,
a ciągłość i gęstość dają $R\widehat J=I$ na
uzupełnieniu. Dowodzi to obu kierunków izomorfizmu.

Dla $\mathcal E$ oznaczmy przez $J_{\mathcal E}$
odwzorowanie kanoniczne na iloczynie algebraicznym.
Dla zwartych $K\Subset\Omega$, $L\Subset\Lambda$
i $m\in\N_0$ mamy
\[
 p_{K\times L,m}(J_{\mathcal E}u)
 \le\pi_{p_{K,m},p_{L,m}}(u),
\]
ponieważ pochodne tensora elementarnego w obu zmiennych
rozdzielają się, a następnie bierzemy infimum
po przedstawieniach. Każdy zwarty podzbiór iloczynu
mieści się w iloczynie swoich zwartych projekcji.
Stąd $J_{\mathcal E}$ jest ciągłe i przedłuża się
do ciągłego $\widehat J_{\mathcal E}$, ponieważ
$\mathcal E(\Omega\times\Lambda)$ jest zupełna.

Wybierzmy przeliczalne lokalnie skończone pokrycia
otwartymi sześcianami $V_i\Subset\Omega$,
$W_j\Subset\Lambda$, oraz gładkie rozkłady jedności
$(\eta_i)$, $(\theta_j)$ podporządkowane tym pokryciom.
Istnienie takich pokryć wynika z wyczerpania
$K_n\subseteq\operatorname{int}K_{n+1}$
(Lemat~\ref{lem:exhaustion17}):
dla warstwy $K_n\setminus\operatorname{int}K_{n-1}$
wybieramy skończenie wiele sześcianów o domknięciach
zawartych w $\operatorname{int}K_{n+1}\setminus K_{n-2}$
(przyjmujemy $K_0=K_{-1}=\varnothing$).
Pokrycie jest lokalnie skończone: dla punktu
w $\operatorname{int}K_N$ to otoczenie nie przecina
sześcianów z warstw $n\ge N+2$, a pozostałych
sześcianów jest skończenie wiele.
Stosujemy Twierdzenie~\ref{thm:dt-smooth-partition}.
Otrzymaną partycję, początkowo indeksowaną przez
rodzinę funkcji przypisanych elementom pokrycia,
grupujemy według tych elementów. Dla ustalonego $V_i$
tylko skończenie wiele niezerowych składników może
być przypisanych do $V_i$, gdyż ich nośniki tworzą
rodzinę lokalnie skończoną, a $\overline V_i$ jest
zwartym podzbiorem $\Omega$. Ich suma $\eta_i$ ma
więc zwarty nośnik zawarty w $V_i$; gdy nie ma
składników, przyjmujemy $\eta_i=0$. Tak samo
konstruujemy $\theta_j$ na $\Lambda$.
Wybieramy odcięcia rzeczywiste
$\chi_i\in\mathcal D(V_i)$, $\psi_j\in\mathcal D(W_j)$,
równe $1$ w otoczeniach tych nośników
(Lemat~\ref{lem:cutoffs17}; przy pustym nośniku
można wybrać odcięcie zerowe).
Istotna jest lokalna skończoność całych rodzin
$(\supp\chi_i)$ i $(\supp\psi_j)$, wynikająca
z lokalnej skończoności pokryć.

Na domknięciach sześcianów stosujemy rzeczywiste
ortonormalne układy trygonometryczne $(e_{ik})_k$,
$(h_{jl})_l$ względem znormalizowanych miar Lebesgue'a.
W każdej współrzędnej bierzemy $1$ oraz
$\sqrt2\cos(\pi n(t-a)/R)$, $\sqrt2\sin(\pi n(t-a)/R)$
dla $n\ge1$, gdzie $a$ jest środkiem, a $2R$ długością
boku. W wielu wymiarach bierzemy iloczyny tych funkcji.
Ortogonalność i normalizacja w jednym wymiarze
wynikają z całkowania iloczynów funkcji sinus
i cosinus po pełnym okresie; w wielu wymiarach
stosujemy twierdzenie Fubiniego. Jest to baza
ortonormalna także nad $\C$, ponieważ jej zespolona
rozpiętość zawiera wykładnicze Fouriera, których
zupełność wykazano w dowodzie
Twierdzenia~\ref{thm:nuclear-cinfty23} (po przesunięciu
środka sześcianu). W przypadku rzeczywistym zupełność
wynika przez rozpatrzenie rzeczywistej funkcji
ortogonalnej do tego układu jako funkcji zespolonej.
Przez $|k|$ i $|l|$ oznaczamy normy euklidesowe
wektorów częstości; każda częstość ma ograniczoną,
zależną tylko od wymiaru, krotność.
Funkcje $e_{ik},h_{jl}$ są funkcjami własnymi
$1-\Delta$ z wartościami własnymi odpowiednio
$1+\pi^2|k|^2/R_i^2$, $1+\pi^2|l|^2/S_j^2$.

Dla $\Phi\in\mathcal E(\Omega\times\Lambda)$
lokalizujemy $\Phi_{ij}=\eta_i\theta_j\Phi$,
przedłużamy zerem, a następnie okresowo z
$\overline V_i\times\overline W_j$.
Gładkość periodyzacji wynika ze znikania w otoczeniu
brzegu. Niech $c_{kl}^{ij}$ będą współczynnikami
w iloczynowej bazie trygonometrycznej: całkami
$\Phi_{ij}e_{ik}h_{jl}$ względem znormalizowanej
miary na iloczynie sześcianów. Kładziemy
$f_{ik}=\chi_i e_{ik}$, $g_{jl}=\psi_j h_{jl}$,
przedłużone zerem poza odpowiednie sześciany.
Dla $a,b,r,s\in\N_0$ całkowanie przez części z
$(1-\Delta_x)^a(1-\Delta_y)^b$ i reguła Leibniza dają
\[
 \begin{aligned}
 |c_{kl}^{ij}|&\le C_{ij,a,b}(1+|k|)^{-2a}
 (1+|l|)^{-2b}
 p_{\overline V_i\times\overline W_j,2a+2b}(\Phi),\\
 p_{K,r}(f_{ik})&\le C_{i,r}(1+|k|)^r,\qquad
 p_{L,s}(g_{jl})\le C_{j,s}(1+|l|)^s.
 \end{aligned}
\]
Stałe są niezależne od $\Phi,k,l$; pochodne ustalonych
odcięć są ograniczone. Dla $2a>r+d_1$, $2b>s+d_2$
szereg iloczynów tych trzech oszacowań jest zbieżny
przez zliczanie punktów kratowych.
Zatem
$T_{ij}(\Phi)=\sum_{k,l}c_{kl}^{ij}f_{ik}\otimes g_{jl}$
istnieje w uzupełnieniu rzutowym i jest liniowe
oraz ciągłe w $\Phi$.
Jednostajna zbieżność szeregu trygonometrycznego
wraz z pochodnymi dowolnego ustalonego rzędu oraz
zupełność układu w $L^2$ identyfikują jego sumę
z periodyzacją $\Phi_{ij}$ na iloczynie sześcianów.
Odcięcia eliminują pozostałe kopie okresowe, więc
$\widehat J_{\mathcal E}T_{ij}(\Phi)=\Phi_{ij}$
na całym $\Omega\times\Lambda$.

Definiujemy $T(\Phi)=\sum_{i,j}T_{ij}(\Phi)$.
W seminormie $\pi_{p_{K,r},p_{L,s}}$ niezerowy wkład
mogą dawać tylko pary z
$K\cap\supp\chi_i\ne\varnothing$ i
$L\cap\supp\psi_j\ne\varnothing$.
Jest ich skończenie wiele. Sieć sum skończonych
jest więc Cauchy'ego w każdej seminormie, a jej
granica istnieje z zupełności.
Dla każdej z tych seminorm $T(\Phi)$ szacujemy
skończoną sumą powyższych oszacowań, co dowodzi
ciągłości $T$. Ciągłość $\widehat J_{\mathcal E}$
oraz lokalnie skończony rozkład jedności dają
$\widehat J_{\mathcal E}T(\Phi)=\sum_{i,j}\Phi_{ij}=\Phi$.

Dla $\Phi(x,y)=f(x)g(y)$ współczynniki $c_{kl}^{ij}$
są iloczynami współczynników $\eta_i f$ i $\theta_j g$,
na mocy Fubiniego. Odpowiednie szeregi w poszczególnych grupach zmiennych
z odcięciami zbiegają w $\mathcal E$
do $\eta_i f$ i $\theta_j g$. Sumy po prostokątnych
zbiorach indeksów $(k,l)$ są tensorami ich sum
częściowych. Ciągłość odwzorowania dwuliniowego
daje $T_{ij}(fg)=(\eta_i f)\otimes(\theta_j g)$.
Sumując po prostokątnych skończonych zbiorach
indeksów $(i,j)$, otrzymujemy tensory sum lokalnych,
które zbiegają w $\mathcal E(\Omega)$ i
$\mathcal E(\Lambda)$ do $f$ i $g$.
Stąd $T J_{\mathcal E}(f\otimes g)=f\otimes g$.
Przez liniowość, ciągłość i gęstość algebraicznego
iloczynu w uzupełnieniu mamy $T\widehat J_{\mathcal E}=I$.
Oba izomorfizmy są zatem udowodnione nad $\R$ i $\C$.
Na mocy Twierdzenia~\ref{thm:nuclear-cinfty23}
i Wniosku~\ref{cor:nuclear-schwartz23} pierwsze
czynniki są nuklearne, więc
Twierdzenie~\ref{thm:nuclear-pi-eps24}
daje również izomorfizmy dla $\varepsilon$.
\end{proof}
